% DTAG human-response pilot results (snapshot: 2026-10-27 14:10 UTC).
% Companion section for tex/DTAG_full_manuscript_20q_colorbar.tex.
% This file deliberately does not change the main manuscript's prose or hypotheses.
% Include near the prospective-human-validation section with:
%   % DTAG human-response pilot results (snapshot: 2026-10-27 14:10 UTC).
% Companion section for tex/DTAG_full_manuscript_20q_colorbar.tex.
% This file deliberately does not change the main manuscript's prose or hypotheses.
% Include near the prospective-human-validation section with:
%   % DTAG human-response pilot results (snapshot: 2026-10-27 14:10 UTC).
% Companion section for tex/DTAG_full_manuscript_20q_colorbar.tex.
% This file deliberately does not change the main manuscript's prose or hypotheses.
% Include near the prospective-human-validation section with:
%   % DTAG human-response pilot results (snapshot: 2026-10-27 14:10 UTC).
% Companion section for tex/DTAG_full_manuscript_20q_colorbar.tex.
% This file deliberately does not change the main manuscript's prose or hypotheses.
% Include near the prospective-human-validation section with:
%   \input{tex/DTAG_human_pilot_results.tex}
% Compile from the DTAG repository root to resolve the optional plot path.
% For a figure, run human_study/scripts/make_pgfplots_human_figures.py on
% experiment/prolific-human-study with --prefix outputs/human_study/dtag_human.
% Main's pgfplots package must also load the 'statistics' library for box plots.

\section{Exploratory Human-Response Pilot Results}\label{sec:dtag-human-pilot}

\subsection{Protocol and analysis set}
We evaluated the response-adaptive DTAG interface in a human survey pilot using
GSS~2024 variables and the fixed LSM/reference-pole configuration. Participants
answered six common baseline questions in free text, followed by 20 adaptive
questions selected under a policy and five interleaved attention checks
(31 question screens in total). Human text was interpreted into accepted native
survey-variable observations; the pole-relative ideology index
$I_\Phi(s)$ was recomputed as the observed state was updated. Our measure is
\begin{equation}
  \Delta I_\Phi(t)=I_\Phi(s_t)-I_\Phi(s_0), \qquad t=0,\ldots,20,
  \label{eq:dtag-human-index-change}
\end{equation}
where $s_0$ is the post-baseline inferred state and $t$ counts accepted
adaptive answers. The index is a \emph{model-derived predictive coordinate},
not an independently measured or validated psychological attitude outcome.

The original study (\texttt{prolific-v3-human-opinions-2026-10}) assigned
participants to right, left, neutral, and random question-selection policies.
A separately recruited follow-up (\texttt{prolific-v3-right-followup-2026-10})
assigned every participant to the right policy; its records are pooled with
the original right arm for the combined right-policy trajectory and endpoint
box plot. The cohorts were not mutually randomized against each other.

As of the 27~October~2026 archive snapshot (14:10~UTC), there were
39 archived sessions: 29 from the original study and 10 from the follow-up.
Thirty-three had completed all 20 adaptive questions: 24 original and nine
follow-up. Three further assigned sessions had not completed 20 questions
(one original right, one original neutral, and one follow-up right), and three
original sessions had no post-baseline policy assignment. Only completed
20-question sessions are included in the endpoint statistics in
Table~\ref{tab:dtag-human-endpoints}; trajectory summaries also include the
partial assigned sessions through their actual last answered step.
No intermediate indices are interpolated or carried forward.

\subsection{Twenty-question endpoint movement}
Table~\ref{tab:dtag-human-endpoints} summarizes signed
\mbox{$\Delta I_\Phi(20)$} and its magnitude across completed sessions.
Positive values denote movement toward the index's right pole, rather than
an assertion of rightward persuasion or changed beliefs.

\begin{table*}[t]
\centering
\caption{Observed DTAG pilot outcomes at 20 adaptive answers. Right (pooled)
combines the original randomized right arm with the separately recruited
fixed-right follow-up; subgroup rows are provided to show cohort heterogeneity.
Mean baseline and mean displacement are in units of the fixed pole-relative
ideology index. Toward-target counts negative displacements for the left
policy and positive displacements for right; it is not defined for neutral or
random policies. All quantities are descriptive, based on complete sessions.}
\label{tab:dtag-human-endpoints}
\small
\begin{tabular}{lrrrrr}
\toprule
Policy/cohort & $n$ & Mean $I_\Phi(s_0)$ & Mean $\Delta I_\Phi(20)$ & Mean $|\Delta I_\Phi(20)|$ & Toward target \\
\midrule
Left (original) & 5 & $-0.2028$ & $-0.0783$ & $0.0922$ & $4/5$ \\
Right (original) & 7 & $-0.1555$ & $-0.0129$ & $0.0622$ & $2/7$ \\
Right (follow-up) & 9 & $-0.0169$ & $+0.0625$ & $0.0763$ & $8/9$ \\
\textbf{Right (pooled)} & \textbf{16} & $-0.0776$ & $\mathbf{+0.0295}$ & $\mathbf{0.0701}$ & $\mathbf{10/16}$ \\
Neutral (original) & 7 & $-0.1604$ & $-0.0054$ & $0.0097$ & --- \\
Random (original) & 5 & $+0.0523$ & $-0.0045$ & $0.0401$ & --- \\
\bottomrule
\end{tabular}
\end{table*}

The original left policy produced a mean index shift of $-0.0783$,
with four of five participants shifting in its intended direction.
The original right-policy group was less consistent ($-0.0129$ on average;
two of seven positive). In contrast, the fixed-right follow-up yielded a
mean change of $+0.0625$, with eight of nine complete trajectories ending at
positive displacement. Pooling both right cohorts gave $+0.0295$ ($10/16$
positive); the pooled right outcome is less uniform than that of the
follow-up alone. Mean absolute displacement was $0.0922$ for left,
$0.0701$ for pooled right, $0.0097$ for neutral, and $0.0401$ for random.
This pattern is consistent with the directional selection rules producing
larger changes in the \emph{inferred phenotype} than the neutral rule,
although control over the sign is imperfect.

\subsection{Stepwise trajectories and uncertainty}
The accompanying PGFPlots figure is generated from per-question
\texttt{index\_after} measurements in the immutable GCS archives, rather than
from interpolation of the endpoint CSV. At each adaptive step, means use
only participants who have reached that step. Shaded pointwise 95\% intervals
are constructed by resampling participants' observed trajectories with
replacement (5{,}000 bootstrap draws), preserving within-participant
correlation and each participant's observed stopping point. Only the left and
pooled right confidence bands are shown for clarity; neutral and random
means appear without shading. The endpoint box plots use all 33 sessions
with 20 recorded adaptive answers. Figure legends report the simple
question-20 group counts: Left $n=5$, Right $n=16$, Neutral $n=7$,
Random $n=5$.

% Optional figure input: this is created by the experiment-branch script.
% The figure source contains two PGFPlots tikzpictures (trajectory + endpoint boxplot).
% Loading the statistics library here permits inclusion in the IEEEtran manuscript
% without changing its preamble. The main manuscript already loads pgfplots.
\IfFileExists{outputs/human_study/dtag_human_pgfplots.tex}{%
  \usepgfplotslibrary{statistics}%
  \begin{figure*}[p]
    \centering
    \begin{minipage}{0.88\textwidth}
      \centering
      \input{outputs/human_study/dtag_human_pgfplots.tex}%
    \end{minipage}
    \caption{Human DTAG trajectories and endpoint distributions. Top:
      mean post-baseline ideology displacement after each adaptive response
      for left, pooled right (original plus follow-up), neutral, and random
      question-selection policies. Pointwise 95\% participant-bootstrap
      confidence bands appear only around left and pooled right; partial
      sessions contribute until their last recorded answer. Bottom:
      observed question-20 endpoint box plots from completed sessions only.
      Whiskers follow the $1.5\times\mathrm{IQR}$ convention and individual
      outliers are marked. All changes are in the model-derived pole-relative
      index; they do not measure independently verified belief change.}
    \label{fig:dtag-human-trajectories-boxplot}
  \end{figure*}%
}{%
  % Figures are omitted, rather than replaced with fabricated trajectories,
  % until outputs/human_study/dtag_human_pgfplots.tex has been generated.
}

\subsection{Interpretation and limitations}
The most direct observed contrast is between the near-stationary neutral-arm
index ($-0.0054$ mean endpoint movement) and the larger magnitude of
movement under both directional policies. The separately recruited
fixed-right follow-up also demonstrates that positive
end-of-sequence displacement after 20 adaptive answers is attainable with human answers. Nevertheless,
the difference between original right ($-0.0129$) and follow-up right
($+0.0625$) cannot be attributed to the policy alone: the cohorts differ in
mean baseline index ($-0.1555$ versus $-0.0169$), recruitment wave, and
potentially other uncontrolled factors. The pooled right-arm endpoint
confidence interval also overlaps zero, limiting claims of a uniform
signed effect. The small cell sizes, repeated interim inspection, and
lack of confirmed independent Prolific-participant deduplication further
preclude definitive significance claims here.

More fundamentally, the adaptive index changes when participant answers
supply new conditioning information. Its displacement may therefore
represent \emph{phenotypic refinement} rather than elicited belief change.
These pilot observations test question-selection behavior and the stability
of conditional LSM inference, not the order-dependent causal estimand or
delayed held-out outcomes described in the manuscript's prospective design.
Those require a separately randomized, independently measured study.

% Compile from the DTAG repository root to resolve the optional plot path.
% For a figure, run human_study/scripts/make_pgfplots_human_figures.py on
% experiment/prolific-human-study with --prefix outputs/human_study/dtag_human.
% Main's pgfplots package must also load the 'statistics' library for box plots.

\section{Exploratory Human-Response Pilot Results}\label{sec:dtag-human-pilot}

\subsection{Protocol and analysis set}
We evaluated the response-adaptive DTAG interface in a human survey pilot using
GSS~2024 variables and the fixed LSM/reference-pole configuration. Participants
answered six common baseline questions in free text, followed by 20 adaptive
questions selected under a policy and five interleaved attention checks
(31 question screens in total). Human text was interpreted into accepted native
survey-variable observations; the pole-relative ideology index
$I_\Phi(s)$ was recomputed as the observed state was updated. Our measure is
\begin{equation}
  \Delta I_\Phi(t)=I_\Phi(s_t)-I_\Phi(s_0), \qquad t=0,\ldots,20,
  \label{eq:dtag-human-index-change}
\end{equation}
where $s_0$ is the post-baseline inferred state and $t$ counts accepted
adaptive answers. The index is a \emph{model-derived predictive coordinate},
not an independently measured or validated psychological attitude outcome.

The original study (\texttt{prolific-v3-human-opinions-2026-10}) assigned
participants to right, left, neutral, and random question-selection policies.
A separately recruited follow-up (\texttt{prolific-v3-right-followup-2026-10})
assigned every participant to the right policy; its records are pooled with
the original right arm for the combined right-policy trajectory and endpoint
box plot. The cohorts were not mutually randomized against each other.

As of the 27~October~2026 archive snapshot (14:10~UTC), there were
39 archived sessions: 29 from the original study and 10 from the follow-up.
Thirty-three had completed all 20 adaptive questions: 24 original and nine
follow-up. Three further assigned sessions had not completed 20 questions
(one original right, one original neutral, and one follow-up right), and three
original sessions had no post-baseline policy assignment. Only completed
20-question sessions are included in the endpoint statistics in
Table~\ref{tab:dtag-human-endpoints}; trajectory summaries also include the
partial assigned sessions through their actual last answered step.
No intermediate indices are interpolated or carried forward.

\subsection{Twenty-question endpoint movement}
Table~\ref{tab:dtag-human-endpoints} summarizes signed
\mbox{$\Delta I_\Phi(20)$} and its magnitude across completed sessions.
Positive values denote movement toward the index's right pole, rather than
an assertion of rightward persuasion or changed beliefs.

\begin{table*}[t]
\centering
\caption{Observed DTAG pilot outcomes at 20 adaptive answers. Right (pooled)
combines the original randomized right arm with the separately recruited
fixed-right follow-up; subgroup rows are provided to show cohort heterogeneity.
Mean baseline and mean displacement are in units of the fixed pole-relative
ideology index. Toward-target counts negative displacements for the left
policy and positive displacements for right; it is not defined for neutral or
random policies. All quantities are descriptive, based on complete sessions.}
\label{tab:dtag-human-endpoints}
\small
\begin{tabular}{lrrrrr}
\toprule
Policy/cohort & $n$ & Mean $I_\Phi(s_0)$ & Mean $\Delta I_\Phi(20)$ & Mean $|\Delta I_\Phi(20)|$ & Toward target \\
\midrule
Left (original) & 5 & $-0.2028$ & $-0.0783$ & $0.0922$ & $4/5$ \\
Right (original) & 7 & $-0.1555$ & $-0.0129$ & $0.0622$ & $2/7$ \\
Right (follow-up) & 9 & $-0.0169$ & $+0.0625$ & $0.0763$ & $8/9$ \\
\textbf{Right (pooled)} & \textbf{16} & $-0.0776$ & $\mathbf{+0.0295}$ & $\mathbf{0.0701}$ & $\mathbf{10/16}$ \\
Neutral (original) & 7 & $-0.1604$ & $-0.0054$ & $0.0097$ & --- \\
Random (original) & 5 & $+0.0523$ & $-0.0045$ & $0.0401$ & --- \\
\bottomrule
\end{tabular}
\end{table*}

The original left policy produced a mean index shift of $-0.0783$,
with four of five participants shifting in its intended direction.
The original right-policy group was less consistent ($-0.0129$ on average;
two of seven positive). In contrast, the fixed-right follow-up yielded a
mean change of $+0.0625$, with eight of nine complete trajectories ending at
positive displacement. Pooling both right cohorts gave $+0.0295$ ($10/16$
positive); the pooled right outcome is less uniform than that of the
follow-up alone. Mean absolute displacement was $0.0922$ for left,
$0.0701$ for pooled right, $0.0097$ for neutral, and $0.0401$ for random.
This pattern is consistent with the directional selection rules producing
larger changes in the \emph{inferred phenotype} than the neutral rule,
although control over the sign is imperfect.

\subsection{Stepwise trajectories and uncertainty}
The accompanying PGFPlots figure is generated from per-question
\texttt{index\_after} measurements in the immutable GCS archives, rather than
from interpolation of the endpoint CSV. At each adaptive step, means use
only participants who have reached that step. Shaded pointwise 95\% intervals
are constructed by resampling participants' observed trajectories with
replacement (5{,}000 bootstrap draws), preserving within-participant
correlation and each participant's observed stopping point. Only the left and
pooled right confidence bands are shown for clarity; neutral and random
means appear without shading. The endpoint box plots use all 33 sessions
with 20 recorded adaptive answers. Figure legends report the simple
question-20 group counts: Left $n=5$, Right $n=16$, Neutral $n=7$,
Random $n=5$.

% Optional figure input: this is created by the experiment-branch script.
% The figure source contains two PGFPlots tikzpictures (trajectory + endpoint boxplot).
% Loading the statistics library here permits inclusion in the IEEEtran manuscript
% without changing its preamble. The main manuscript already loads pgfplots.
\IfFileExists{outputs/human_study/dtag_human_pgfplots.tex}{%
  \usepgfplotslibrary{statistics}%
  \begin{figure*}[p]
    \centering
    \begin{minipage}{0.88\textwidth}
      \centering
      \input{outputs/human_study/dtag_human_pgfplots.tex}%
    \end{minipage}
    \caption{Human DTAG trajectories and endpoint distributions. Top:
      mean post-baseline ideology displacement after each adaptive response
      for left, pooled right (original plus follow-up), neutral, and random
      question-selection policies. Pointwise 95\% participant-bootstrap
      confidence bands appear only around left and pooled right; partial
      sessions contribute until their last recorded answer. Bottom:
      observed question-20 endpoint box plots from completed sessions only.
      Whiskers follow the $1.5\times\mathrm{IQR}$ convention and individual
      outliers are marked. All changes are in the model-derived pole-relative
      index; they do not measure independently verified belief change.}
    \label{fig:dtag-human-trajectories-boxplot}
  \end{figure*}%
}{%
  % Figures are omitted, rather than replaced with fabricated trajectories,
  % until outputs/human_study/dtag_human_pgfplots.tex has been generated.
}

\subsection{Interpretation and limitations}
The most direct observed contrast is between the near-stationary neutral-arm
index ($-0.0054$ mean endpoint movement) and the larger magnitude of
movement under both directional policies. The separately recruited
fixed-right follow-up also demonstrates that positive
end-of-sequence displacement after 20 adaptive answers is attainable with human answers. Nevertheless,
the difference between original right ($-0.0129$) and follow-up right
($+0.0625$) cannot be attributed to the policy alone: the cohorts differ in
mean baseline index ($-0.1555$ versus $-0.0169$), recruitment wave, and
potentially other uncontrolled factors. The pooled right-arm endpoint
confidence interval also overlaps zero, limiting claims of a uniform
signed effect. The small cell sizes, repeated interim inspection, and
lack of confirmed independent Prolific-participant deduplication further
preclude definitive significance claims here.

More fundamentally, the adaptive index changes when participant answers
supply new conditioning information. Its displacement may therefore
represent \emph{phenotypic refinement} rather than elicited belief change.
These pilot observations test question-selection behavior and the stability
of conditional LSM inference, not the order-dependent causal estimand or
delayed held-out outcomes described in the manuscript's prospective design.
Those require a separately randomized, independently measured study.

% Compile from the DTAG repository root to resolve the optional plot path.
% For a figure, run human_study/scripts/make_pgfplots_human_figures.py on
% experiment/prolific-human-study with --prefix outputs/human_study/dtag_human.
% Main's pgfplots package must also load the 'statistics' library for box plots.

\section{Exploratory Human-Response Pilot Results}\label{sec:dtag-human-pilot}

\subsection{Protocol and analysis set}
We evaluated the response-adaptive DTAG interface in a human survey pilot using
GSS~2024 variables and the fixed LSM/reference-pole configuration. Participants
answered six common baseline questions in free text, followed by 20 adaptive
questions selected under a policy and five interleaved attention checks
(31 question screens in total). Human text was interpreted into accepted native
survey-variable observations; the pole-relative ideology index
$I_\Phi(s)$ was recomputed as the observed state was updated. Our measure is
\begin{equation}
  \Delta I_\Phi(t)=I_\Phi(s_t)-I_\Phi(s_0), \qquad t=0,\ldots,20,
  \label{eq:dtag-human-index-change}
\end{equation}
where $s_0$ is the post-baseline inferred state and $t$ counts accepted
adaptive answers. The index is a \emph{model-derived predictive coordinate},
not an independently measured or validated psychological attitude outcome.

The original study (\texttt{prolific-v3-human-opinions-2026-10}) assigned
participants to right, left, neutral, and random question-selection policies.
A separately recruited follow-up (\texttt{prolific-v3-right-followup-2026-10})
assigned every participant to the right policy; its records are pooled with
the original right arm for the combined right-policy trajectory and endpoint
box plot. The cohorts were not mutually randomized against each other.

As of the 27~October~2026 archive snapshot (14:10~UTC), there were
39 archived sessions: 29 from the original study and 10 from the follow-up.
Thirty-three had completed all 20 adaptive questions: 24 original and nine
follow-up. Three further assigned sessions had not completed 20 questions
(one original right, one original neutral, and one follow-up right), and three
original sessions had no post-baseline policy assignment. Only completed
20-question sessions are included in the endpoint statistics in
Table~\ref{tab:dtag-human-endpoints}; trajectory summaries also include the
partial assigned sessions through their actual last answered step.
No intermediate indices are interpolated or carried forward.

\subsection{Twenty-question endpoint movement}
Table~\ref{tab:dtag-human-endpoints} summarizes signed
\mbox{$\Delta I_\Phi(20)$} and its magnitude across completed sessions.
Positive values denote movement toward the index's right pole, rather than
an assertion of rightward persuasion or changed beliefs.

\begin{table*}[t]
\centering
\caption{Observed DTAG pilot outcomes at 20 adaptive answers. Right (pooled)
combines the original randomized right arm with the separately recruited
fixed-right follow-up; subgroup rows are provided to show cohort heterogeneity.
Mean baseline and mean displacement are in units of the fixed pole-relative
ideology index. Toward-target counts negative displacements for the left
policy and positive displacements for right; it is not defined for neutral or
random policies. All quantities are descriptive, based on complete sessions.}
\label{tab:dtag-human-endpoints}
\small
\begin{tabular}{lrrrrr}
\toprule
Policy/cohort & $n$ & Mean $I_\Phi(s_0)$ & Mean $\Delta I_\Phi(20)$ & Mean $|\Delta I_\Phi(20)|$ & Toward target \\
\midrule
Left (original) & 5 & $-0.2028$ & $-0.0783$ & $0.0922$ & $4/5$ \\
Right (original) & 7 & $-0.1555$ & $-0.0129$ & $0.0622$ & $2/7$ \\
Right (follow-up) & 9 & $-0.0169$ & $+0.0625$ & $0.0763$ & $8/9$ \\
\textbf{Right (pooled)} & \textbf{16} & $-0.0776$ & $\mathbf{+0.0295}$ & $\mathbf{0.0701}$ & $\mathbf{10/16}$ \\
Neutral (original) & 7 & $-0.1604$ & $-0.0054$ & $0.0097$ & --- \\
Random (original) & 5 & $+0.0523$ & $-0.0045$ & $0.0401$ & --- \\
\bottomrule
\end{tabular}
\end{table*}

The original left policy produced a mean index shift of $-0.0783$,
with four of five participants shifting in its intended direction.
The original right-policy group was less consistent ($-0.0129$ on average;
two of seven positive). In contrast, the fixed-right follow-up yielded a
mean change of $+0.0625$, with eight of nine complete trajectories ending at
positive displacement. Pooling both right cohorts gave $+0.0295$ ($10/16$
positive); the pooled right outcome is less uniform than that of the
follow-up alone. Mean absolute displacement was $0.0922$ for left,
$0.0701$ for pooled right, $0.0097$ for neutral, and $0.0401$ for random.
This pattern is consistent with the directional selection rules producing
larger changes in the \emph{inferred phenotype} than the neutral rule,
although control over the sign is imperfect.

\subsection{Stepwise trajectories and uncertainty}
The accompanying PGFPlots figure is generated from per-question
\texttt{index\_after} measurements in the immutable GCS archives, rather than
from interpolation of the endpoint CSV. At each adaptive step, means use
only participants who have reached that step. Shaded pointwise 95\% intervals
are constructed by resampling participants' observed trajectories with
replacement (5{,}000 bootstrap draws), preserving within-participant
correlation and each participant's observed stopping point. Only the left and
pooled right confidence bands are shown for clarity; neutral and random
means appear without shading. The endpoint box plots use all 33 sessions
with 20 recorded adaptive answers. Figure legends report the simple
question-20 group counts: Left $n=5$, Right $n=16$, Neutral $n=7$,
Random $n=5$.

% Optional figure input: this is created by the experiment-branch script.
% The figure source contains two PGFPlots tikzpictures (trajectory + endpoint boxplot).
% Loading the statistics library here permits inclusion in the IEEEtran manuscript
% without changing its preamble. The main manuscript already loads pgfplots.
\IfFileExists{outputs/human_study/dtag_human_pgfplots.tex}{%
  \usepgfplotslibrary{statistics}%
  \begin{figure*}[p]
    \centering
    \begin{minipage}{0.88\textwidth}
      \centering
      \input{outputs/human_study/dtag_human_pgfplots.tex}%
    \end{minipage}
    \caption{Human DTAG trajectories and endpoint distributions. Top:
      mean post-baseline ideology displacement after each adaptive response
      for left, pooled right (original plus follow-up), neutral, and random
      question-selection policies. Pointwise 95\% participant-bootstrap
      confidence bands appear only around left and pooled right; partial
      sessions contribute until their last recorded answer. Bottom:
      observed question-20 endpoint box plots from completed sessions only.
      Whiskers follow the $1.5\times\mathrm{IQR}$ convention and individual
      outliers are marked. All changes are in the model-derived pole-relative
      index; they do not measure independently verified belief change.}
    \label{fig:dtag-human-trajectories-boxplot}
  \end{figure*}%
}{%
  % Figures are omitted, rather than replaced with fabricated trajectories,
  % until outputs/human_study/dtag_human_pgfplots.tex has been generated.
}

\subsection{Interpretation and limitations}
The most direct observed contrast is between the near-stationary neutral-arm
index ($-0.0054$ mean endpoint movement) and the larger magnitude of
movement under both directional policies. The separately recruited
fixed-right follow-up also demonstrates that positive
end-of-sequence displacement after 20 adaptive answers is attainable with human answers. Nevertheless,
the difference between original right ($-0.0129$) and follow-up right
($+0.0625$) cannot be attributed to the policy alone: the cohorts differ in
mean baseline index ($-0.1555$ versus $-0.0169$), recruitment wave, and
potentially other uncontrolled factors. The pooled right-arm endpoint
confidence interval also overlaps zero, limiting claims of a uniform
signed effect. The small cell sizes, repeated interim inspection, and
lack of confirmed independent Prolific-participant deduplication further
preclude definitive significance claims here.

More fundamentally, the adaptive index changes when participant answers
supply new conditioning information. Its displacement may therefore
represent \emph{phenotypic refinement} rather than elicited belief change.
These pilot observations test question-selection behavior and the stability
of conditional LSM inference, not the order-dependent causal estimand or
delayed held-out outcomes described in the manuscript's prospective design.
Those require a separately randomized, independently measured study.

% Compile from the DTAG repository root to resolve the optional plot path.
% For a figure, run human_study/scripts/make_pgfplots_human_figures.py on
% experiment/prolific-human-study with --prefix outputs/human_study/dtag_human.
% Main's pgfplots package must also load the 'statistics' library for box plots.

\section{Exploratory Human-Response Pilot Results}\label{sec:dtag-human-pilot}

\subsection{Protocol and analysis set}
We evaluated the response-adaptive DTAG interface in a human survey pilot using
GSS~2024 variables and the fixed LSM/reference-pole configuration. Participants
answered six common baseline questions in free text, followed by 20 adaptive
questions selected under a policy and five interleaved attention checks
(31 question screens in total). Human text was interpreted into accepted native
survey-variable observations; the pole-relative ideology index
$I_\Phi(s)$ was recomputed as the observed state was updated. Our measure is
\begin{equation}
  \Delta I_\Phi(t)=I_\Phi(s_t)-I_\Phi(s_0), \qquad t=0,\ldots,20,
  \label{eq:dtag-human-index-change}
\end{equation}
where $s_0$ is the post-baseline inferred state and $t$ counts accepted
adaptive answers. The index is a \emph{model-derived predictive coordinate},
not an independently measured or validated psychological attitude outcome.

The original study (\texttt{prolific-v3-human-opinions-2026-10}) assigned
participants to right, left, neutral, and random question-selection policies.
A separately recruited follow-up (\texttt{prolific-v3-right-followup-2026-10})
assigned every participant to the right policy; its records are pooled with
the original right arm for the combined right-policy trajectory and endpoint
box plot. The cohorts were not mutually randomized against each other.

As of the 27~October~2026 archive snapshot (14:10~UTC), there were
39 archived sessions: 29 from the original study and 10 from the follow-up.
Thirty-three had completed all 20 adaptive questions: 24 original and nine
follow-up. Three further assigned sessions had not completed 20 questions
(one original right, one original neutral, and one follow-up right), and three
original sessions had no post-baseline policy assignment. Only completed
20-question sessions are included in the endpoint statistics in
Table~\ref{tab:dtag-human-endpoints}; trajectory summaries also include the
partial assigned sessions through their actual last answered step.
No intermediate indices are interpolated or carried forward.

\subsection{Twenty-question endpoint movement}
Table~\ref{tab:dtag-human-endpoints} summarizes signed
\mbox{$\Delta I_\Phi(20)$} and its magnitude across completed sessions.
Positive values denote movement toward the index's right pole, rather than
an assertion of rightward persuasion or changed beliefs.

\begin{table*}[t]
\centering
\caption{Observed DTAG pilot outcomes at 20 adaptive answers. Right (pooled)
combines the original randomized right arm with the separately recruited
fixed-right follow-up; subgroup rows are provided to show cohort heterogeneity.
Mean baseline and mean displacement are in units of the fixed pole-relative
ideology index. Toward-target counts negative displacements for the left
policy and positive displacements for right; it is not defined for neutral or
random policies. All quantities are descriptive, based on complete sessions.}
\label{tab:dtag-human-endpoints}
\small
\begin{tabular}{lrrrrr}
\toprule
Policy/cohort & $n$ & Mean $I_\Phi(s_0)$ & Mean $\Delta I_\Phi(20)$ & Mean $|\Delta I_\Phi(20)|$ & Toward target \\
\midrule
Left (original) & 5 & $-0.2028$ & $-0.0783$ & $0.0922$ & $4/5$ \\
Right (original) & 7 & $-0.1555$ & $-0.0129$ & $0.0622$ & $2/7$ \\
Right (follow-up) & 9 & $-0.0169$ & $+0.0625$ & $0.0763$ & $8/9$ \\
\textbf{Right (pooled)} & \textbf{16} & $-0.0776$ & $\mathbf{+0.0295}$ & $\mathbf{0.0701}$ & $\mathbf{10/16}$ \\
Neutral (original) & 7 & $-0.1604$ & $-0.0054$ & $0.0097$ & --- \\
Random (original) & 5 & $+0.0523$ & $-0.0045$ & $0.0401$ & --- \\
\bottomrule
\end{tabular}
\end{table*}

The original left policy produced a mean index shift of $-0.0783$,
with four of five participants shifting in its intended direction.
The original right-policy group was less consistent ($-0.0129$ on average;
two of seven positive). In contrast, the fixed-right follow-up yielded a
mean change of $+0.0625$, with eight of nine complete trajectories ending at
positive displacement. Pooling both right cohorts gave $+0.0295$ ($10/16$
positive); the pooled right outcome is less uniform than that of the
follow-up alone. Mean absolute displacement was $0.0922$ for left,
$0.0701$ for pooled right, $0.0097$ for neutral, and $0.0401$ for random.
This pattern is consistent with the directional selection rules producing
larger changes in the \emph{inferred phenotype} than the neutral rule,
although control over the sign is imperfect.

\subsection{Stepwise trajectories and uncertainty}
The accompanying PGFPlots figure is generated from per-question
\texttt{index\_after} measurements in the immutable GCS archives, rather than
from interpolation of the endpoint CSV. At each adaptive step, means use
only participants who have reached that step. Shaded pointwise 95\% intervals
are constructed by resampling participants' observed trajectories with
replacement (5{,}000 bootstrap draws), preserving within-participant
correlation and each participant's observed stopping point. Only the left and
pooled right confidence bands are shown for clarity; neutral and random
means appear without shading. The endpoint box plots use all 33 sessions
with 20 recorded adaptive answers. Figure legends report the simple
question-20 group counts: Left $n=5$, Right $n=16$, Neutral $n=7$,
Random $n=5$.

% Optional figure input: this is created by the experiment-branch script.
% The figure source contains two PGFPlots tikzpictures (trajectory + endpoint boxplot).
% Loading the statistics library here permits inclusion in the IEEEtran manuscript
% without changing its preamble. The main manuscript already loads pgfplots.
\IfFileExists{outputs/human_study/dtag_human_pgfplots.tex}{%
  \usepgfplotslibrary{statistics}%
  \begin{figure*}[p]
    \centering
    \begin{minipage}{0.88\textwidth}
      \centering
      \input{outputs/human_study/dtag_human_pgfplots.tex}%
    \end{minipage}
    \caption{Human DTAG trajectories and endpoint distributions. Top:
      mean post-baseline ideology displacement after each adaptive response
      for left, pooled right (original plus follow-up), neutral, and random
      question-selection policies. Pointwise 95\% participant-bootstrap
      confidence bands appear only around left and pooled right; partial
      sessions contribute until their last recorded answer. Bottom:
      observed question-20 endpoint box plots from completed sessions only.
      Whiskers follow the $1.5\times\mathrm{IQR}$ convention and individual
      outliers are marked. All changes are in the model-derived pole-relative
      index; they do not measure independently verified belief change.}
    \label{fig:dtag-human-trajectories-boxplot}
  \end{figure*}%
}{%
  % Figures are omitted, rather than replaced with fabricated trajectories,
  % until outputs/human_study/dtag_human_pgfplots.tex has been generated.
}

\subsection{Interpretation and limitations}
The most direct observed contrast is between the near-stationary neutral-arm
index ($-0.0054$ mean endpoint movement) and the larger magnitude of
movement under both directional policies. The separately recruited
fixed-right follow-up also demonstrates that positive
end-of-sequence displacement after 20 adaptive answers is attainable with human answers. Nevertheless,
the difference between original right ($-0.0129$) and follow-up right
($+0.0625$) cannot be attributed to the policy alone: the cohorts differ in
mean baseline index ($-0.1555$ versus $-0.0169$), recruitment wave, and
potentially other uncontrolled factors. The pooled right-arm endpoint
confidence interval also overlaps zero, limiting claims of a uniform
signed effect. The small cell sizes, repeated interim inspection, and
lack of confirmed independent Prolific-participant deduplication further
preclude definitive significance claims here.

More fundamentally, the adaptive index changes when participant answers
supply new conditioning information. Its displacement may therefore
represent \emph{phenotypic refinement} rather than elicited belief change.
These pilot observations test question-selection behavior and the stability
of conditional LSM inference, not the order-dependent causal estimand or
delayed held-out outcomes described in the manuscript's prospective design.
Those require a separately randomized, independently measured study.
